% SPDX-License-Identifier: Apache-2.0
%
% PCIe in TxHDL. Built by //docs:pcie. Every listing is included from
% the tree, and every printout and figure is what the build made by
% running the example.
\documentclass[journal,9pt]{IEEEtran}

% Listings of Rust set by rustfmt at 80 columns: 5.6pt is what fits
% the frame of a column (issue 1061).
\newcommand{\listingsize}{\fontsize{5.6}{6.6}\selectfont}
% SPDX-License-Identifier: Apache-2.0
%
% The house preamble, shared by every document in this directory. Loaded
% after \documentclass, before \title. Keeping it in one file is what
% stops two articles drifting apart in font, listing style or figure
% style.
% Computer Modern. IEEEtran selects Times (ptm) by default, so the face has
% to be chosen explicitly.
%
% Latin Modern is the Computer Modern variety used here, and the choice is
% forced rather than aesthetic. Under T1 encoding, plain `cmr` has no Type 1
% outlines unless cm-super is installed, and it is not in the pinned
% distribution, so pdflatex falls back to EC bitmaps and every glyph in the
% PDF becomes a Type 3 font. That was measured: the first build produced a
% document with 10 Type 3 fonts and no Type 1. Latin Modern is Computer
% Modern redrawn with a full T1 repertoire and real .pfb outlines, so the
% same design comes out scalable.
\usepackage[T1]{fontenc}
\usepackage[utf8]{inputenc}
\usepackage{lmodern}
% microtype is not in the pinned TeX distribution. emergencystretch alone
% keeps the two-column measure from producing overfull lines.
\setlength{\emergencystretch}{3em}

\usepackage{amsmath}
\usepackage{amssymb}
\usepackage{amsthm}
\usepackage{booktabs}
\usepackage{array}
% enumitem is not in the pinned TeX distribution; the list spacing below
% does what \setlist would have done.
\usepackage{float}
\usepackage{xcolor}
\usepackage{listings}
\usepackage{tikz}
\usetikzlibrary{arrows.meta,positioning,fit,backgrounds,calc}
\usepackage[hidelinks,bookmarks=true]{hyperref}

% Numbered, definition-style box for a rule the reader is meant to apply.
\theoremstyle{definition}
\newtheorem{rulebox}{Rule}
\newtheorem{example}{Example}

% Unnumbered asides.
\theoremstyle{remark}
\newtheorem*{tip}{Tip}
\newtheorem*{pitfall}{Pitfall}

% Tighter lists, in place of enumitem's \setlist.
\setlength{\topsep}{2pt}
\setlength{\itemsep}{1pt}
\setlength{\parsep}{0pt}

% Code in the running text. The typewriter face under T1 ligates `<<`
% and `>>` into guillemets, so a shift printed as \code{<<} came out as
% a French quotation mark (issue 571). microtype's \DisableLigatures is
% not in the pinned distribution, but the primitive it rests on is
% pdfTeX's own: \pdfnoligatures strips every ligature from the font it
% is given, and the current font inside \texttt is the typewriter face
% at the size in use, so each size is stripped the first time \code
% sets it. Code wants no ligature at all: two quotes are two quotes.
\newcommand{\code}[1]{\texttt{\ifdefined\pdfnoligatures\pdfnoligatures\font\fi #1}}
\newcommand{\term}[1]{\emph{#1}}

% Syntax highlighting. Four roles, distinguishable in colour and still
% distinguishable in greyscale, because an IEEE article gets printed: the
% keyword is the only bold role, the comment is the only italic one, and the
% two remaining roles differ in lightness as well as in hue.
\definecolor{kwblue}{RGB}{20,60,150}
\definecolor{cmtgreen}{RGB}{90,120,90}
\definecolor{strred}{RGB}{150,50,40}
\definecolor{numgray}{gray}{0.45}
\definecolor{framegray}{gray}{0.70}
\definecolor{bgshade}{gray}{0.97}
% A darker fill for figure cells. The listing background has to stay very
% light to keep code readable; a figure cell has to be clearly shaded or the
% distinction it draws is invisible in print.
\definecolor{cellshade}{gray}{0.90}

% The listing size is the document's choice, because it depends on the
% column: a two-column article sets `\listingsize` to 6pt and keeps its
% listings within 70 columns, or to 5.6pt when it includes source that
% rustfmt sets at 80, which is what fits a column's frame (issue 1061);
% a one-column document keeps the body size and includes source files
% at 80. `//docs:frame_test` checks every listing line against its frame. Lines are never broken; a line that does not fit
% overflows the frame, which is visible, rather than wrapping, which is
% not.
\providecommand{\listingsize}{\scriptsize}

% `'` is deliberately not a string delimiter: the language uses it for loop
% labels, as Rust does, so treating it as a quote would swallow the rest of
% the line.
\lstdefinestyle{txhdl}{
  basicstyle=\ttfamily\listingsize,
  keywordstyle=\bfseries\color{kwblue},
  commentstyle=\itshape\color{cmtgreen},
  stringstyle=\color{strred},
  numberstyle=\tiny\color{numgray},
  morekeywords={mod,use,pub,const,type,struct,enum,trait,impl,fn,async,let,mut,
    match,if,else,loop,while,for,in,break,continue,return,as,await,
    module,bus,channel,transaction,protocol,pipeline,flow,seq,config,
    port,stage,pipe,instance,connect,bind,tag,domain,blackbox,foreign,
    property,role,signal,handler,true,false,
    sequence,interface,modport,implements,package,import,var,wire,func,proc,
    case,default,next,fork,branch,join,state,fsm},
  morecomment=[l]{//},
  morestring=[b]",
  showstringspaces=false,
  columns=fullflexible,
  keepspaces=true,
  breaklines=false,
  % Framed, so a listing is bounded rather than floating in the column.
  frame=single,
  rulecolor=\color{framegray},
  framesep=3pt,
  backgroundcolor=\color{bgshade},
  xleftmargin=3pt,
  xrightmargin=3pt,
  aboveskip=6pt,
  belowskip=6pt,
  % Every listing is numbered and captioned, so it can be referred to by
  % number. `caption` without `float` numbers the listing and leaves it
  % where it was written, which is what a listing in running prose needs.
  captionpos=b,
  abovecaptionskip=2pt,
  belowcaptionskip=6pt,
}

% Figures drawn as text. Framed the same way, and with the highlighting
% switched off, because a box-and-arrow diagram is not code and half its
% words turning blue reads as an accident. Setting a style adds keys rather
% than clearing the ones already set, so the three roles are neutralised by
% hand rather than by leaving them out. Program output is printed in this
% style too, and a netlist's assignment can run past the frame, so a line
% that does not fit is broken and the rest indented; source never is.
\lstdefinestyle{ascii}{
  basicstyle=\ttfamily\listingsize,
  keywordstyle=\ttfamily,
  commentstyle=\ttfamily,
  stringstyle=\ttfamily,
  columns=fullflexible,
  keepspaces=true,
  breaklines=true,
  breakindent=2em,
  frame=single,
  rulecolor=\color{framegray},
  framesep=3pt,
  backgroundcolor=\color{bgshade},
  xleftmargin=3pt,
  xrightmargin=3pt,
  aboveskip=6pt,
  belowskip=6pt,
}
\lstset{style=txhdl}

% House TikZ styles. `shorten` lives on the arrow styles rather than on each
% arrow, so every arrow in the document gets the gap, including ones added
% later. TikZ clips a path at a node's border, so without it an arrowhead
% butts against the box with no gap at all. Keep `node distance` well above
% twice the shorten, or the arrow shrinks to nothing.
\tikzset{
  box/.style={draw,rounded corners=2pt,align=center,inner sep=4pt,
              font=\footnotesize},
  boxf/.style={box,fill=bgshade},
  lbl/.style={font=\scriptsize\itshape,align=center},
  fl/.style={-{Stealth[length=4pt]},shorten >=2.5pt,shorten <=2.5pt},
  fld/.style={fl,dashed},
}


% The contents list: the class gives a section's number 2.75em, which
% a Roman numeral past thirty overruns onto the title, so the box is
% wider here, and a subsection's entry starts where a title does.
\makeatletter
\def\l@section#1#2{\addpenalty{\@secpenalty}\addvspace{1.0em plus 1pt}%
    \@tempdima 5.75em \begingroup \parindent \z@ \rightskip \@pnumwidth%
    \parfillskip-\@pnumwidth {\bfseries\leavevmode #1}\hfil\hbox to\@pnumwidth{\hss #2}\par%
    \endgroup}
\def\l@subsection{\@dottedtocline{2}{5.75em}{5em}}
\def\l@subsubsection{\@dottedtocline{3}{10.75em}{5.5em}}
\makeatother

% SPDX-License-Identifier: Apache-2.0
% The boards' memory maps, written by //tools/memmap from the maps the
% designs decode (issue 444). A document asks for a table with
% \mmget{<what>}{<board>}, and a table the generator does not write
% stops the document rather than leaving a gap.
\input{tools/memmap/mm_tables}
\makeatletter
\newcommand{\mmget}[2]{%
  \@ifundefined{mm@#1@#2}%
    {\PackageError{memmap}{//tools/memmap writes no #1 for #2}%
      {Add it to tools/memmap/main.rs.}}%
    {\csname mm@#1@#2\endcsname}}
\makeatother


\InputIfFileExists{buildstamp.tex}{}{}
\providecommand{\buildstamp}{Draft build (outside the Bazel flow).}

% A range of a source file, by its begin/end markers.
\newcommand{\pinspart}[3]{%
\lstinputlisting[rangebeginprefix=//\ begin\{,rangebeginsuffix=\},%
rangeendprefix=//\ end\{,rangeendsuffix=\},includerangemarker=false,%
linerange=#1-#1,caption={#2},label={#3}]{lib/parts/src/bus/axi_pins.rs}}
\newcommand{\perpinspart}[3]{%
\lstinputlisting[rangebeginprefix=//\ begin\{,rangebeginsuffix=\},%
rangeendprefix=//\ end\{,rangeendsuffix=\},includerangemarker=false,%
linerange=#1-#1,caption={#2},label={#3}]{lib/parts/src/bus/axi_per_pins.rs}}
\newcommand{\barpart}[3]{%
\lstinputlisting[rangebeginprefix=//\ begin\{,rangebeginsuffix=\},%
rangeendprefix=//\ end\{,rangeendsuffix=\},includerangemarker=false,%
linerange=#1-#1,caption={#2},label={#3}]{pcie/src/bar.rs}}

\title{PCIe in TxHDL}

\author{Filip Filmar and Dragiša Janković%
\thanks{Both authors are independent. Filip Filmar is the
corresponding author, at f@hdlfactory.com; Dragiša Janković is
reached at linkedin.com/in/dragisa-jankovic.}%
\thanks{This tutorial details integrating TxHDL logic behind a PCIe endpoint,
building upon the AXI bus architecture~\cite{axi}. All code listings are
embedded from the project repository \code{HDL/txhdl} during compilation, and
all traces reflect automated testbench execution. The project overview~\cite{cover}
indexes all companion reports.}%
\thanks{The authors used Claude, an AI assistant, while developing
architectural concepts, documentation, and software. Verbatim prompts are
preserved in repository commit logs.}%
\thanks{\buildstamp}}

\markboth{PCIe in TxHDL}{Filmar and Janković: PCIe in TxHDL}

\begin{document}
\maketitle

\begin{abstract}
The Alinx AX7A200B FPGA board includes a PCIe $\times$2 edge connector wired to
dedicated GTP transceiver quads and integrated PCIe hard blocks. AMD supplies
the physical and data link layers as the DMA/Bridge Subsystem for PCI Express IP
core. The build system configures this core for Gen2 operation across two lanes
with its AXI-Bridge bypass window enabled, mapping TxHDL logic behind the PCIe
interface. Inbound reads and writes targeting BAR1 emerge from the core as AXI4
transactions at external master pins. A dedicated library component,
\code{AxiPins}, converts between raw pin-level signaling and TxHDL channel
interfaces, connecting to a memory-mapped register unit containing four 64-bit
registers. The pipeline is validated via cycle-accurate co-simulation in
\code{nvc} and Verilator, exercised through high-level Rust drivers, and
implemented through Vivado synthesis and placement to produce a timing-clean
bitstream.
\end{abstract}

\section{What Is Here}
\label{sec:p-what}

This design spans three functional modules within the repository tree:

\emph{Bus interface adapter.} Component \code{AxiPins}, located in
\code{txhdl\_parts} under \code{bus::axi\_pins}, bridges external pin-level AXI4
master signals to internal TxHDL channels. It drops in wherever a host tracker
attaches to a physical bus interface.

\emph{Target register file.} Module \code{pcie::bar}, implemented under
\code{//pcie}, services transactions mapped to BAR1. Unit \code{BarRegs} provides
four 64-bit registers behind a peripheral bus tracker. Composite unit
\code{PcieBar} integrates \code{AxiPins}, the tracker, and the register file
into a single lowered component.

\emph{PCIe endpoint and board integration.} Rule \code{//pcie:xdma\_x2}
generates the vendor PCIe controller via Vivado, while \code{pcie/board}
provides the top-level structural wrapper and physical constraints.

\section{The Endpoint}
\label{sec:p-endpoint}

The Bazel build system executes Vivado toolchains hermetically. Extension rules
in \code{rules\_vivado} generate vendor IP cores through \code{vivado\_ip}, which
takes core parameters and produces HDL sources alongside the \code{.xci}
descriptor. Downstream synthesis rules import this descriptor via
\code{read\_ip}, compiling vendor logic together with application RTL.

Two AMD cores provide AXI master interfaces mapped to PCIe Base Address
Registers on Artix-7 devices: the AXI Memory-Mapped PCIe bridge
(\code{axi\_pcie}) and the DMA/Bridge Subsystem for PCI Express (\code{xdma}).
Both were evaluated at Gen2 link speeds across two physical lanes using board pin
assignments. Both mapped cleanly to the two GTP transceiver channels and the
integrated PCIe hard macro. However, \code{axi\_pcie} failed timing closure by
0.097\,ns across three internal paths within its unused AXI slave bridge, which
cannot be disabled. Conversely, \code{xdma} allows disabling unused slave ports,
closing timing with positive slack. At Gen2 $\times$2, the core requires a
125\,MHz user clock frequency, rejecting 62.5\,MHz configurations.

Configuration parameters reside in \code{pcie/BUILD.bazel}: Gen2 negotiation,
two physical lanes, a 100\,MHz PCIe reference clock, a 64-bit master bus width,
and an active bypass aperture mapping a 1~MiB address window to BAR1. On Artix-7
devices, the XDMA core does not support pure bridge mode; setting bridge mode
silently defaults to a DMA engine configuration where BAR0 maps internal DMA
registers and drives a descriptor-based master~\cite{i451}. The bypass window
resolves this limitation by providing direct bridging: host transactions
targeting BAR1 emerge directly on the \code{m\_axib} master port without
descriptor intervention. Outbound master addresses span 64 bits, of which
downstream logic decodes the lower 32 bits.

\section{The Pins}
\label{sec:p-pins}

External IP blocks acting as AXI4 masters expose standard pin interfaces:
individual \code{valid} and \code{ready} control lines per channel and discrete
wires for data and address payloads. In contrast, TxHDL components communicate
across typed channels transporting structured transactions~\cite{axi}.
Component \code{AxiPins} bridges these two conventions with zero pipeline latency.

\pinspart{ports}{The pins, as the unit's two sides.}{lst:p-ports}

Listing~\ref{lst:p-ports} illustrates \code{AxiPins} port declarations. The 24
master interface pins aggregate into a structured type, \code{AxiHostPins}, which
derives \code{Ports}. Input interfaces hold these pins as a grouped field,
\code{pins}, alongside inbound response channels. During lowering, individual
signals map to standard top-level ports such as \code{inp\_pins\_awid}, while
outbound transactions drive \code{outp\_aw}. Designs incorporating host pins
encapsulate them into composite structs, mirroring the approach used for Vreteno
JTAG interfaces (issue~\#579).

When an upstream master asserts an address phase or write beat, \code{AxiPins}
presents it to the corresponding channel whenever downstream capacity is
available, asserting the host's \code{ready} signal. Conversely, write responses
and read beats arriving at channel heads drive host \code{valid} and data lines,
retiring whenever the host asserts \code{ready}. Macro \code{select!} converts
between raw AXI4 wire encodings and strongly-typed TxHDL enums. Upstream masters
omit \code{AxQOS} and \code{AxREGION} signaling, so the adapter drives both
outputs to zero.

\pinspart{unit}{\code{AxiPins}: the step.}{lst:p-unit}

Listing~\ref{lst:p-unit} details the synchronous step logic of \code{AxiPins}.
Because pins model zero-delay combinatorial wires, driving processes evaluate in
the same simulation delta cycle as consuming units, matching physical netlist
behaviors.

Example driver \code{ex\_axi\_pins} evaluates this adapter in simulation
(Listing~\ref{lst:p-ex}). A behavioral host model drives
\code{AxiPins<32, 64, 8, 4>}, matching endpoint bus widths, connected to a
peripheral tracker and target RAM. The test issues three 64-bit write beats,
reads back four words, and attempts an out-of-bounds read that triggers a bus
fault. Listing~\ref{lst:p-out} records simulation output, and
Figure~\ref{fig:p-pins} diagrams transaction handshakes across the pin boundary.
Lowered netlists match software execution under both \code{nvc} and Verilator
regressions: \code{//docs:axi\_pins\_sim\_axi\_pins\_tb\_test} and
\code{//docs:axi\_pins\_vsim\_test}.

\lstinputlisting[caption={\code{ex\_axi\_pins.rs}: a host on pins, and a
RAM behind them.},label={lst:p-ex}]{lib/examples/ex_axi_pins.rs}

\lstinputlisting[style=ascii,caption={What \code{ex\_axi\_pins} printed
when the build ran it.},label={lst:p-out}]{out_axi_pins.txt}

\begin{figure}[htbp]
\centering
\input{axi_pins_timing.tex}
\caption{The write through the pins: its address phase and its three
beats handed to the channels as the host offers them, and its response
handed back.}
\label{fig:p-pins}
\end{figure}

\section{The Pins the Other Way}
\label{sec:p-perpins}

External IP blocks can also function as AXI4 peripherals, such as the AMD DDR3
memory controller evaluated in issue~\#1023. Component \code{AxiPerPins} provides
the inverse adapter: it attaches to peripheral channel endpoints where a
peripheral tracker typically resides, exposing physical slave pins without
internal buffering.

\perpinspart{ports}{The pins, the other way.}{lst:p-perports}

Listing~\ref{lst:p-perports} shows port definitions for \code{AxiPerPins}.
Peripheral-driven signals comprise eleven pins, including handshake lines and
response data, encapsulated within struct \code{AxiPerDriven}. The adapter
presents address phases, write beats, and response acknowledge lines as outputs
alongside channel returns to the system link.

\perpinspart{unit}{\code{AxiPerPins}: the step.}{lst:p-perunit}

Listing~\ref{lst:p-perunit} presents the synchronous step logic. Address phases
and write beats at channel heads assert external \code{valid} lines and data
fields whenever the peripheral asserts \code{ready}. Inbound write responses and
read data transfer to link channels when channel buffer space is available. The
adapter exports all address phase fields except region qualifiers, which
single-port memory controllers do not inspect. Quality-of-service bits
(\code{AxQOS}) are forwarded directly.

Example testbench \code{ex\_axi\_per\_pins} verifies this interface by attaching
a simulated memory model behind \code{AxiPerPins<30, 32, 4, 5>}, matching DDR3
controller widths. Listing~\ref{lst:p-perex} shows testbench construction,
Listing~\ref{lst:p-perout} lists output traces, and Figure~\ref{fig:p-perpins}
charts write transaction timing across peripheral pins. Multi-simulator test
suites under \code{nvc} and Verilator confirm formal equivalence with generated
netlists: \code{//docs:axi\_per\_pins\_sim\_axi\_per\_pins\_tb\_test} and
\code{//docs:axi\_per\_pins\_vsim\_test}.

\lstinputlisting[caption={\code{ex\_axi\_per\_pins.rs}: a memory on
pins, and a client in front.},label={lst:p-perex}]{lib/examples/ex_axi_per_pins.rs}

\lstinputlisting[style=ascii,caption={What \code{ex\_axi\_per\_pins}
printed when the build ran it.},label={lst:p-perout}]{out_axi_per_pins.txt}

\begin{figure}[htbp]
\centering
\input{axi_per_pins_timing.tex}
\caption{The write the other way: its address phase and its three
beats taken off the link's channels as the peripheral, always ready,
takes them, and its response sent back on the link.}
\label{fig:p-perpins}
\end{figure}

\section{Behind BAR1}
\label{sec:p-bar}

\code{BarRegs} is four words, at the offsets a host reads and writes in
BAR1.

\mmget{regs}{bar}

The identifier reads \mmget{ident}{bar}, which encodes \code{TxHDL} in ASCII
along with revision metadata. The build generates this table directly from
\code{WORDS} declarations in \code{bar.rs}, and test cases access each register
at its declared offset.

The register bank services single transactions sequentially, mirroring Vreteno
system timer semantics. Read requests complete in the cycle of acceptance, and
writes commit upon arrival of the write data beat. Decoding logic inspects
address bits 4 and 3, relying on upstream trackers to filter out-of-range
requests. Writes update complete 64-bit words regardless of byte enable strobes.

\barpart{regs}{\code{BarRegs}: four words behind a peripheral
tracker.}{lst:p-regs}

Composite module \code{PcieBar} integrates the three components into a single
unit: the \code{AxiPins} adapter, peripheral tracker, and register bank, joined
by internal channels. External ports expose AXI master pins, system resets, and
diagnostic LED outputs, connecting directly to top-level XDMA interfaces.

\barpart{bar}{\code{PcieBar}: everything behind BAR1, as one
unit.}{lst:p-bar}

Diagnostic test suite \code{//pcie:bar\_test} exercises \code{PcieBar} across
its pin interface, validating identifier queries, scratchpad read-write loops,
LED state toggles, and write counter tracking. Automated checks confirm that the
synthesized netlist forms a single module with pins exposed as top-level ports.

\section{The Board}
\label{sec:p-board}

Top-level integration module \code{pcie\_top.v} coordinates components beyond the
scope of lowered units: differential reference clock buffers, the XDMA endpoint
macro, and clock-reset distribution networks. \code{PcieBar} executes
synchronously on the 125\,MHz PCIe user clock, avoiding clock domain crossing
logic. Active-low resets drive lowered logic through synchronous inversion.

\lstinputlisting[caption={\code{pcie/board/pcie\_top.v}: the
endpoint and the design behind it.},label={lst:p-top}]{pcie/board/pcie_top.v}

Physical constraints in \code{pcie/board/ax7a200\_pcie.xdc} assign the two
serial lanes, 100\,MHz differential reference clock inputs, reset pushbuttons,
and diagnostic LEDs. High-speed gigabit transceivers omit standard CMOS I/O
declarations. Because the AX7A200B hardware manual omits PCIe slot reset (PERST\#)
routing~\cite{i179}, system resets map to the on-board pushbutton.

\lstinputlisting[caption={\code{pcie/board/ax7a200\_pcie.xdc}: the
board's pins.},label={lst:p-xdc}]{pcie/board/ax7a200_pcie.xdc}

Compilation via \code{bazel build //pcie:endpoint\_synth} synthesizes the XDMA
IP, lowers \code{PcieBar}, and executes Vivado synthesis without critical
warnings. Implementation rule \code{//pcie:endpoint\_pnr} places and routes the
design in eight minutes. The design meets all timing constraints, achieving
0.656\,ns setup slack and 0.036\,ns hold slack across 42\,813 endpoints on the
100\,MHz reference clock and derived 125\,MHz user clock. Device utilization
consumes 11\,320 slice LUTs, 13\,570 flip-flops, 20.5 Block RAM tiles, two GTP
transceivers, and the integrated PCIe block.

\section{What Is Not Done}
\label{sec:p-not}

Hardware validation within an active host motherboard remains to be completed.
Physical host testing will verify Gen2 link negotiation across both lanes, PCIe
device enumeration, Base Address Register mapping, and register read-write
cycles from host software~\cite{i180}.

The endpoint is thoroughly verified in behavioral simulation. Compilation
scripts build simulation libraries directly from Vivado simulation export assets.
Testbench \code{//pcie/sim:xdma\_test} evaluates top-level gateware and
\code{BarRegs} without external stimuli, confirming the link remains quiescent.
Testbench \code{//pcie/sim:rp\_test} connects the design to a behavioral PCIe
root port model provided by AMD example designs~\cite{i178}. The simulation
establishes a Gen2 $\times$2 link, interrogates vendor and device IDs, programs
BAR1, and executes register validation test suite \code{pcie/sim/rp/tests.vh}.
The test sequence reads identifier strings, performs scratchpad read-write
loops, and verifies write operation counters, passing all checks in 26 minutes
of simulation. Initial runs without the bypass aperture exposed DMA engine
defaults at BAR0~\cite{i451}.

Rule \code{vivado\_ip} masks certain IP configuration errors, requiring manual
inspection of synthesis logs when parameters change~\cite{i177}. Furthermore,
PCIe slot reset line PERST\# remains unrouted on this hardware
revision~\cite{i179}. Currently, \code{PcieBar} operates on an isolated bus
clocked by the PCIe endpoint. Bridging the BAR aperture to the primary Vreteno
processor bus will require dual-clock asynchronous FIFOs, implemented using
\code{ChanCdc} channel synchronizers~\cite{i131}.

\IEEEtriggeratref{2}
\begin{thebibliography}{9}

\bibitem{axi}
F.~Filmar and D.~Janković, ``AXI in TxHDL,'' project repository
\code{HDL/txhdl}, \code{//docs:axi}, 2026.

\bibitem{i131}
Issue 131, ``feat(test): co-simulate a unit of several clocks, and
lower a clock-crossing channel,'' project repository \code{HDL/txhdl},
2026.

\bibitem{i177}
Issue 177, ``fix(rules\_vivado): vivado\_ip succeeds when Vivado
rejects the IP's configuration,'' project repository
\code{HDL/txhdl}, 2026.

\bibitem{i178}
Issue 178, ``feat(pcie): simulate the PCIe endpoint, its IP and a root
port,'' project repository \code{HDL/txhdl}, 2026.

\bibitem{i179}
Issue 179, ``board: the AX7A200B's PCIe reset, PERST\#, and its pin,''
project repository \code{HDL/txhdl}, 2026.

\bibitem{i180}
Issue 180, ``test(pcie): run the endpoint in a PC and read its BAR from
the host,'' project repository \code{HDL/txhdl}, 2026.

\bibitem{i451}
Issue 451, ``fix(pcie): the endpoint is a DMA engine, not a bridge;
the Artix-7 XDMA ignores functional\_mode,'' project repository
\code{HDL/txhdl}, 2026.

\bibitem{cover}
F.~Filmar and D.~Janković, ``TxHDL documents,'' project repository
\code{HDL/txhdl}, \code{//docs:cover}, 2026.

\end{thebibliography}

\end{document}
